\documentclass{article}
\usepackage{enumitem}
\usepackage{amsmath}
\begin{document}

\title{Chapter 17: Breathing and Exchange of Gases}
\date{}
\maketitle

\begin{enumerate}[label=Q\arabic*.]

\item The partial pressure of oxygen (pO$_2$) in alveolar air is approximately:
\\(A) 159 mmHg
\\(B) 104 mmHg
\\(C) 40 mmHg
\\(D) 760 mmHg

\item Oxygen is primarily transported in blood as:
\\(A) Dissolved in plasma (about 97\%)
\\(B) Bound to haemoglobin in RBCs (about 97\%)
\\(C) Bound to albumin
\\(D) As bicarbonate ions

\item The Bohr effect describes:
\\(A) Increased O$_2$ affinity of haemoglobin at high CO$_2$/low pH
\\(B) Decreased O$_2$ affinity of haemoglobin at high CO$_2$/low pH (right shift of ODC)
\\(C) Transport of CO$_2$ as carbaminohaemoglobin
\\(D) Effect of temperature on lung volume

\item The chloride shift (Hamburger phenomenon) in RBCs involves:
\\(A) Cl$^-$ moving out of RBC as HCO$_3^-$ moves in
\\(B) HCO$_3^-$ moving out of RBC as Cl$^-$ moves in, maintaining electrical neutrality
\\(C) Na$^+$ and Cl$^-$ exchanging across the RBC membrane
\\(D) CO$_2$ converting to HCO$_3^-$ inside plasma

\item Assertion: CO$_2$ is transported primarily as bicarbonate ions in blood.\\
Reason: CO$_2$ reacts with water inside RBCs (catalyzed by carbonic anhydrase) to form carbonic acid, which dissociates to H$^+$ and HCO$_3^-$.
\\(A) Both true, Reason is correct explanation
\\(B) Both true, Reason is not correct explanation
\\(C) Assertion true, Reason false
\\(D) Assertion false, Reason true

\item The tidal volume (TV) is:
\\(A) Maximum air that can be forcibly expelled after normal expiration
\\(B) Volume of air inhaled or exhaled during normal quiet breathing ($\approx$500 mL)
\\(C) Air remaining in lungs after maximum expiration
\\(D) Maximum air lungs can hold

\item Residual volume (RV) cannot be expelled because:
\\(A) Diaphragm cannot contract further
\\(B) Air trapped in non-collapsible alveoli prevents complete lung collapse
\\(C) Bronchioles close during forced expiration
\\(D) Intercostal muscles limit chest compression

\item The respiratory center in the brain that controls the basic rhythm of breathing is located in the:
\\(A) Cerebral cortex
\\(B) Medulla oblongata
\\(C) Hypothalamus
\\(D) Pons only

\item Surfactant in alveoli is produced by:
\\(A) Type I pneumocytes
\\(B) Type II pneumocytes (great alveolar cells)
\\(C) Alveolar macrophages
\\(D) Ciliated epithelial cells

\item In COPD (emphysema), the total lung capacity increases but vital capacity decreases because:
\\(A) Airways are blocked by mucus
\\(B) Alveolar walls are destroyed, air trapping increases RV while gas exchange surface decreases
\\(C) Bronchial constriction increases residual volume
\\(D) Diaphragm muscles become hyperactive

\item Haemoglobin has a higher affinity for CO compared to O$_2$ by approximately:
\\(A) 10 times
\\(B) 100 times
\\(C) 250 times
\\(D) 1000 times

\item The vital capacity (VC) of lungs is:
\\(A) TV + IRV + ERV
\\(B) TV + IRV + ERV + RV
\\(C) TV + RV
\\(D) IRV + ERV

\item Which of the following changes stimulates increase in breathing rate?
\\(A) Increased pO$_2$ in blood
\\(B) Increased pCO$_2$ and decreased pH in blood
\\(C) Decreased pCO$_2$ in blood
\\(D) Increased pH of cerebrospinal fluid

\item Match the following lung volumes/capacities:\\
1. IRV $\rightarrow$ (a) Extra air inspired after normal inspiration ($\approx$2500 mL)\\
2. ERV $\rightarrow$ (b) Extra air expired after normal expiration ($\approx$1000 mL)\\
3. RV $\rightarrow$ (c) Air remaining after maximum expiration ($\approx$1100 mL)\\
4. FRC $\rightarrow$ (d) ERV + RV
\\(A) 1-a, 2-b, 3-c, 4-d
\\(B) 1-b, 2-a, 3-d, 4-c
\\(C) 1-c, 2-d, 3-a, 4-b
\\(D) 1-d, 2-c, 3-b, 4-a

\item Pneumothorax (air in pleural cavity) causes lung collapse because:
\\(A) Excess air in alveoli increases pressure
\\(B) The negative intrapleural pressure that keeps lungs inflated is lost
\\(C) Diaphragm is paralyzed
\\(D) Surfactant is destroyed by air entry

\end{enumerate}
\end{document}
